Mit $C_v = \frac{f}{2}R$ und $C_p = C_v + R = \frac{f+2}{2}R$ (FoTaBe S.~171) ist
\begin{align}
 \kappa &= \frac{C_p}{C_v} = \frac{f+2}{f} = 1+\frac{2}{f}
 \qquad\Leftrightarrow\qquad
 \result{f = \frac{2}{\kappa-1}}
\end{align}

\emph{Test} mit den Adiabatenexponenten aus der Tabelle \emph{Thermische Daten von Gasen} (FoTaBe S.~191):
\begin{align}
 \text{Argon:}\quad f &= \frac{2}{\kArO-1} = \fAr \\
 \text{Helium:}\quad f &= \frac{2}{\kHeO-1} = \fHe \\
 \text{Stickstoff:}\quad f &= \frac{2}{\kNzO-1} = \fNz \\
 \text{Sauerstoff:}\quad f &= \frac{2}{\kOzO-1} = \fOz
\end{align}
Die Edelgase haben wie erwartet etwa drei Freiheitsgrade (nur Translation), die zweiatomigen Gase etwa fünf (Translation und zwei Rotationen).

\emph{Methan und Freon 12:}
\begin{align}
 \text{Methan:}\quad f &= \frac{2}{\kMeO-1} = \fMe \approx \result{\fMeP{0}{2}} \\
 \text{Freon 12:}\quad f &= \frac{2}{\kFrO-1} = \fFr \approx \result{\fFrP{0}{2}}
\end{align}
Methan (CH$_4$) ist ein räumliches Molekül mit drei Translations- und drei Rotationsfreiheitsgraden; der Rest stammt von teilweise angeregten Schwingungen. Beim schweren Molekül Freon 12 (CCl$_2$F$_2$) sind viele Schwingungen schon bei Zimmertemperatur angeregt, deshalb ist $f$ viel grösser.
